\documentclass[a4paper,11pt]{article}

\usepackage[utf8]{inputenc}
\usepackage[T1]{fontenc}
\usepackage[ngerman]{babel}
\usepackage{amsmath}
\usepackage{amssymb}
\usepackage{xcolor}
\usepackage{colortbl}
\usepackage{array}
\usepackage{tikz}
\usepackage{needspace}
\usepackage{geometry}
\geometry{a4paper,margin=2.2cm}

% ----- Dark Mode (Pflicht) -----
\definecolor{bgdark}{HTML}{1E1E1E}
\definecolor{fgtext}{HTML}{E6E6E6}
\definecolor{accent}{HTML}{89B4FA}
\definecolor{accent2}{HTML}{F9E2AF}
\definecolor{muted}{HTML}{A6ADC8}
\definecolor{gridcol}{HTML}{4A4A4A}
\definecolor{rulecol}{HTML}{6A6A6A}
\pagecolor{bgdark}
\color{fgtext}

\setlength{\arrayrulewidth}{0.3pt}
\arrayrulecolor{rulecol}
\setlength{\parindent}{0pt}
\setlength{\parskip}{4pt}

% ----- Karo-Raster zum Ausfüllen -----
\newcommand{\karo}[1]{\par\noindent
  \begin{tikzpicture}
    \draw[step=5mm, gridcol, very thin] (0,0) grid (\linewidth,#1);
  \end{tikzpicture}\par}

% ----- Laufende Kopfzeile -----
\usepackage{fancyhdr}
\pagestyle{fancy}
\fancyhf{}
\renewcommand{\headrulewidth}{0pt}
\fancyhead[L]{\footnotesize\color{muted}\textsc{Numerische Methoden} $\cdot$ Pr\"ufung 03 Probeklausur}
\fancyhead[R]{\footnotesize\color{muted}Matr.-Nr.\ \textcolor{rulecol}{\rule{2.4cm}{0.3pt}}\quad \thepage}

% ----- Aufgabenkopf -----
\newcommand{\aufg}[3]{%
  \par\vspace{2pt}
  \noindent\textcolor{rulecol}{\rule{\linewidth}{0.3pt}}\\[3pt]
  {\large\bfseries\color{accent} Aufgabe #1\ \ \textcolor{fgtext}{\normalsize\mdseries #2}}\hfill
  {\bfseries\color{accent2}#3 Punkte}\\[1pt]
  \noindent\textcolor{rulecol}{\rule{\linewidth}{0.3pt}}
  \par\vspace{6pt}
}
\newcommand{\tp}[1]{\hfill{\small\color{muted}\textit{#1 Punkte}}}

\begin{document}

% =====================================================================
% DECKBLATT
% =====================================================================
\thispagestyle{fancy}

{\small\color{muted}\textsc{DHBW Ravensburg} \\ Duale Hochschule Baden-W\"urttemberg}

\vspace{10pt}
{\fontsize{34}{38}\selectfont\bfseries\color{accent} Pr\"ufung}\\[2pt]
{\itshape\color{fgtext} Numerische Methoden $\cdot$ Pr\"ufung 03 Probeklausur}

\vspace{10pt}
\noindent\textcolor{rulecol}{\rule{\linewidth}{0.4pt}}
\vspace{4pt}

\renewcommand{\arraystretch}{1.32}
\noindent\begin{tabular}{@{}p{4.2cm}p{11cm}@{}}
\textsc{\color{muted}Studiengang}     & W4DSKI\_108 Data Science und K\"unstliche Intelligenz\\
\textsc{\color{muted}Jahrgang}        & RV-WDS125 / RV-WDS225\\
\textsc{\color{muted}Modul}           & Fortgeschrittene Lineare Algebra und Analysis\\
\textsc{\color{muted}Veranstaltung}   & Numerische Methoden\\
\textsc{\color{muted}Pr\"ufer}        & Prof.\ Dr.-Ing.\ Mark Schutera\\
\textsc{\color{muted}Datum}           & \rule{4cm}{0.3pt}\\
\textsc{\color{muted}Bearbeitungszeit}& 90 Minuten\\
\textsc{\color{muted}Max.\ Punktzahl} & 80 Punkte\\
\textsc{\color{muted}Hilfsmittel}     & kein Taschenrechner. Open book; alle gedruckten oder
                                        handschriftlichen Unterlagen, die auf Papier mitgebracht werden.\\
\end{tabular}

\vspace{6pt}
\noindent\textcolor{rulecol}{\rule{\linewidth}{0.4pt}}

\vspace{10pt}
{\small\color{muted}\textsc{Anweisungen}}\par
\vspace{2pt}
{\small
Schreiben Sie Ihre L\"osungen und L\"osungswege leserlich in die vorgesehenen Bereiche.
Ordnen Sie jede L\"osung deutlich der zugeh\"origen Aufgabe zu. \textbf{Begr\"unden Sie Ihre
Antworten kurz: eine reine Zahl ohne Rechenweg gibt keine volle Punktzahl.} Br\"uche d\"urfen
als L\"osung stehen bleiben; Ableitungen d\"urfen analytisch gebildet werden, sofern nicht
ausdr\"ucklich anders verlangt. Lesen Sie die Aufgabenstellung sorgf\"altig, bevor Sie mit der
Bearbeitung beginnen. Viel Erfolg!}

\vspace{8pt}
{\small\color{muted}\textsc{Themenabdeckung.} Diese Probeklausur deckt ausschlie\ss lich den Stoff der
Vorlesungsnotizen \glqq Numerical Methods\grqq{} ab (Kapitel 1--7): Grundlagen \& Gleitkomma,
Fehleranalyse, Newton/Sekante, globale Optimierung, numerische Integration, Modelltraining.}

\vspace{12pt}
{\small\color{muted}\textsc{Durch Pr\"ufer auszuf\"ullen}}\par
\vspace{3pt}
\renewcommand{\arraystretch}{1.4}
\noindent\begin{tabular}{|l|c|c|c|c|c|c|c|}
\hline
\color{muted}Aufgabe & 1 & 2 & 3 & 4 & 5 & 6 & $\Sigma$\\
\hline
\color{muted}Maximal & 18 & 10 & 14 & 14 & 14 & 10 & 80\\
\hline
\color{muted}Erreicht & & & & & & & \\
\hline
\end{tabular}

% =====================================================================
% AUFGABE 1 -- Kapitel 1-2
% =====================================================================
\clearpage
\aufg{1}{Gleitkomma, Festkomma, Diskretisierung \& Brute Force}{18}

\textbf{(a)} Sie entwerfen ein Gleitkommasystem (floating-point system) mit $1$ Vorzeichenbit
und $4$ weiteren Bits, die Sie frei auf Mantisse (mantissa) und Exponent (exponent) verteilen.
\tp{5}

\smallskip
Diskutieren Sie den Trade-off: Welche Eigenschaft des Systems verbessert sich, wenn Sie ein Bit
der \emph{Mantisse} zuweisen, welche, wenn Sie es dem \emph{Exponenten} geben? Argumentieren Sie
\"uber Aufl\"osung (resolution) und Wertebereich (dynamic range).

\smallskip
Fixieren wir die Aufteilung auf $2$ Mantissenbits und $2$ Exponentenbits mit Bias $1$.
Bestimmen Sie das Maschinenepsilon $\varepsilon_{\mathrm{mach}}$ und berechnen Sie in diesem
System
\[
   \frac{1}{\bigl(1+\tfrac{1}{10}\bigr)-1}.
\]
Geben Sie den Nenner nach Rundung sowie das Endergebnis an, vergleichen Sie mit dem
mathematischen Wert und erkl\"aren Sie, warum hier die Ausl\"oschung (catastrophic cancellation)
das Resultat zerst\"ort.

\karo{7.5cm}

\clearpage
\textbf{(b)} Ein \emph{Festkommasystem} (fixed-point) verwende $8$ Bit, davon $4$ Vorkomma- und
$4$ Nachkommabits (Skalierungsfaktor $2^{-4}$).
\tp{4}

\smallskip
Geben Sie die Aufl\"osung (kleinste darstellbare Differenz) und die gr\"o\ss te darstellbare
(vorzeichenlose) Zahl an. Stellen Sie $5{,}25$ in diesem System dar (als ganzzahliger Speicherwert
und bin\"ar). Vergleichen Sie Festkomma und Gleitkomma: Wie unterscheiden sich die Abst\"ande
benachbarter darstellbarer Zahlen \"uber den Wertebereich?

\karo{6.5cm}

\clearpage
\textbf{(c)} Wir diskretisieren das Intervall $[0,3]$ mit $N=6$ gleich gro\ss en Teilintervallen.
\tp{4}

\smallskip
Berechnen Sie die Schrittweite $h=\frac{b-a}{N}$ und geben Sie alle Gitterpunkte $x_i$ an.
Ordnen Sie anschlie\ss end die folgenden drei Fehlerquellen jeweils einem Fehlertyp zu
(\emph{Modellfehler}, \emph{Abschneidefehler}, \emph{Rundungsfehler}) und begr\"unden Sie kurz:
\begin{enumerate}\setlength{\itemsep}{1pt}
\item[(i)]   Die Sigmoid-Funktion $\sigma(x)=\tfrac{1}{1+e^{-x}}$ wird durch ihre Taylor-Reihe
             nach dem quadratischen Glied ersetzt.
\item[(ii)]  Ein Bremsweg wird ohne Ber\"ucksichtigung des Reifenprofils modelliert.
\item[(iii)] Das Ergebnis $\tfrac{10}{3}$ wird als $3{,}333$ gespeichert.
\end{enumerate}

\karo{6cm}

\clearpage
\textbf{(d)} Wir suchen die L\"osung des linearen Systems
\[
  \begin{pmatrix} 1 & 2 \\ 3 & 1 \end{pmatrix}
  \begin{pmatrix} w \\ b \end{pmatrix}
  =
  \begin{pmatrix} 4 \\ 5 \end{pmatrix}
\]
mit einer Gitter-Suche (grid search) \"uber $w,b\in\{0,1,2\}$.
\tp{5}

\smallskip
W\"ahlen Sie ein Fehlerma\ss\ (error) und begr\"unden Sie Ihre Wahl kurz. Bestimmen Sie damit
$w^\ast$ und $b^\ast$. Vergleichen Sie das Ergebnis mit der exakten L\"osung. Erkl\"art die
Gitter-Suche das System exakt? Wie k\"onnten Sie die Methode verbessern -- und ist es
grunds\"atzlich m\"oglich, so eine exakte L\"osung zu finden?

\smallskip
{\small\color{muted}\textsc{Hinweis.} Die exakte L\"osung lautet $w^\star=6/5$ und $b^\star=7/5$.}

\karo{8cm}

% =====================================================================
% AUFGABE 2 -- Kapitel 3
% =====================================================================
\clearpage
\aufg{2}{Fehleranalyse: die vier Eigenschaften}{10}

Wir suchen das Minimum der Funktion $g(x)=(x-2)^2$ auf $[0,4]$ mit st\"uckweise konstanter
gitterbasierter Diskretisierung bei zwei Schrittweiten $h_1=1$ und $h_2=\tfrac12$. Die
diskretisierte Funktion $\hat g$ wird am \emph{n\"achstgelegenen} Gitterpunkt ausgewertet. Die
Eingabe sei durch einen Messfehler um $\tilde x=x\pm\tfrac12$ gest\"ort.

\begin{center}
\begin{tikzpicture}[scale=0.85]
  \draw[->,rulecol] (-0.3,0) -- (4.6,0) node[right]{\small$x$};
  \draw[->,rulecol] (0,-0.3) -- (0,4.6) node[above]{\small$g(x)$};
  \foreach \x in {1,2,3,4} \draw[rulecol] (\x,-0.08) -- (\x,0.08) node[below=2pt]{\scriptsize$\x$};
  \foreach \y in {2,4} \draw[rulecol] (-0.08,\y) -- (0.08,\y) node[left=2pt]{\scriptsize$\y$};
  \draw[accent,thick,domain=0:4,samples=80,smooth] plot (\x,{(\x-2)*(\x-2)});
\end{tikzpicture}
\end{center}

Analysieren Sie nacheinander die vier Eigenschaften der Fehleranalyse. \emph{Notation:} $g$ exakt,
$\hat g$ Verfahren, $\tilde x$ gest\"orte Eingabe.
\tp{10}

\smallskip
\emph{(i) Kondition (conditioning):} Quantifizieren Sie die Kondition $\kappa=|g(x)-g(\tilde x)|$
des Problems an mindestens \emph{zwei} Stellen Ihrer Wahl im Intervall und zeigen Sie so den
Unterschied zwischen gut und schlecht konditioniertem Bereich.

\smallskip
\emph{(ii) Stabilit\"at (stability):} Berechnen Sie die Stabilit\"at
$s=\frac{|\hat g(\tilde x)-\hat g(x)|}{|\tilde x-x|}$ der diskretisierten Funktion an einem
Referenzpunkt Ihrer Wahl. Welche Schrittweite f\"uhrt zur stabileren Auswertung? Begr\"unden Sie.

\smallskip
\emph{(iii) Konsistenz (consistency):} Quantifizieren Sie den Konsistenzfehler
$c=|\hat g(x)-g(x)|$ f\"ur $h_1$ und $h_2$ an einer Stelle, die nicht auf dem groben Gitter liegt.
Welche Schrittweite ist konsistenter?

\smallskip
\emph{(iv) Konvergenz (convergence):} Verbinden Sie Stabilit\"at und Konsistenz. F\"uhrt eine
kleinere Schrittweite bei numerischen Verfahren \emph{zwingend} zu einer besseren Approximation?
Begr\"unden Sie kurz.

\karo{8.5cm}

% =====================================================================
% AUFGABE 3 -- Kapitel 4
% =====================================================================
\clearpage
\aufg{3}{Newton-, Taylor- \& Sekantenverfahren}{14}

\textbf{(a)} Wir suchen eine Nullstelle (root) von $f(x)=x^2-2$ (also $\sqrt 2$).
\tp{8}

\begin{center}
\begin{tikzpicture}[scale=0.85]
  \draw[->,rulecol] (-0.3,0) -- (2.8,0) node[right]{\small$x$};
  \draw[->,rulecol] (0,-2.4) -- (0,3.0) node[above]{\small$f(x)$};
  \foreach \x in {1,2} \draw[rulecol] (\x,-0.08) -- (\x,0.08) node[below=2pt]{\scriptsize$\x$};
  \foreach \y in {-2,2} \draw[rulecol] (-0.08,\y) -- (0.08,\y) node[left=2pt]{\scriptsize$\y$};
  \draw[accent,thick,domain=-0.2:2.45,samples=80,smooth] plot (\x,{(\x)*(\x)-2});
  \draw[accent2,thick] (2,2) circle (2.5pt);
  \node[accent2] at (2.35,2.4) {\scriptsize$x_0$};
\end{tikzpicture}
\end{center}

\textbf{Leiten Sie} das Newton-Verfahren $x_{n+1}=x_n-\frac{f(x_n)}{f'(x_n)}$ \emph{aus der
Taylor-Entwicklung erster Ordnung} von $f$ um $x_n$ her. F\"uhren Sie dann mit $x_0=2$ zwei
Iterationen durch ($x_1,x_2$ d\"urfen als Bruch angegeben werden) und nennen Sie das
Konvergenzkriterium f\"ur das Newton-Verfahren. Was bedeutet \emph{quadratische Konvergenz}
praktisch?

\karo{8.5cm}

\clearpage
\textbf{(b)} Nehmen Sie nun an, $f$ sei nicht in geschlossener Form gegeben, sondern nur durch
Funktionsauswertungen zug\"anglich.
\tp{6}

\smallskip
Geben Sie eine numerische Approximation (Vorw\"artsdifferenz) f\"ur $f'(x_0)$ an und beziffern Sie
die \emph{Ordnung} des Approximationsfehlers. Leiten Sie daraus das \emph{Sekantenverfahren} ab und
geben Sie seine Iterationsformel an. Welchen Versagensmodus der Newton-Nullstellensuche umgehen Sie
damit -- und welchen \emph{nicht}? F\"uhren Sie schlie\ss lich \emph{einen} Sekanten-Schritt mit
$x_0=2,\,x_1=1{,}5$ f\"ur $f(x)=x^2-2$ aus.

\karo{8.5cm}

% =====================================================================
% AUFGABE 4 -- Kapitel 5
% =====================================================================
\clearpage
\aufg{4}{Globale Optimierung}{14}

\textbf{(a)} Erkl\"aren Sie kurz den Unterschied zwischen einem \emph{lokalen} und einem
\emph{globalen} Minimum. Warum kann eine lokale Methode (z.\,B.\ Newton) auf einer multimodalen
Landschaft wie den \emph{Shekel-Foxholes}
$f(x)=-\sum_i \frac{c_i}{(x-a_i)^2+r_i}$ das globale Minimum verfehlen? Was bedeutet
\emph{Konvexit\"at} f\"ur dieses Problem?
\tp{4}

\karo{5cm}

\clearpage
\textbf{(b)} F\"ur die Funktion $g(x)=x^3-3x$ f\"uhren Sie \emph{einen} Schritt des
Newton-Verfahrens \emph{f\"ur Optima}
\[
   x_{n+1}=x_n-\frac{g'(x_n)}{\lvert g''(x_n)\rvert}
\]
mit Startwert $x_0=2$ durch. Klassifizieren Sie das angesteuerte Extremum mit dem
Test der zweiten Ableitung. Wozu dient der Betrag $\lvert g''\rvert$ im Nenner?
\tp{5}

\karo{7cm}

\clearpage
\textbf{(c)} Bei \emph{zuf\"alligen Neustarts} (random restarts) treffe ein einzelner Start mit
Wahrscheinlichkeit $p=\tfrac12$ das Becken (basin) des globalen Minimums.
\tp{3}

\smallskip
Berechnen Sie die Erfolgswahrscheinlichkeit $P=1-(1-p)^K$ f\"ur $K=3$ Neustarts und bestimmen
Sie das kleinste $K$ mit $P>0{,}9$.

\karo{5cm}

\clearpage
\textbf{(d)} Wir wechseln zur Minimumsuche auf einer \emph{rauen} Verlustlandschaft mit einem
hochfrequenten St\"orterm,
\[
   f(x)=(x-1)^2+\tfrac12\sin(4\pi x),\qquad x\in[0,2].
\]
Erkl\"aren Sie kurz, warum der hochfrequente Sinus-Term f\"ur lokale Optimierer problematisch ist.
Nennen Sie neben der globalen Optimierung einen weiteren Ansatz, um diese Problematik zu umgehen,
und w\"ahlen Sie eine daf\"ur geeignete Schrittweite $h$.
\tp{2}

\karo{5cm}

% =====================================================================
% AUFGABE 5 -- Kapitel 6
% =====================================================================
\clearpage
\aufg{5}{Numerische Integration \& Interpolation}{14}

\textbf{(a)} Geben Sie f\"ur ein Intervall $[a,b]$ mit $h=b-a$ und Mittelpunkt $m=\tfrac{a+b}{2}$
die Formeln der Quadraturregeln \emph{Rechteck} (links/rechts), \emph{Mittelpunkt}, \emph{Trapez}
und \emph{Simpson} an. Ordnen Sie sie nach steigender Genauigkeit und nennen Sie die jeweilige
akkumulierte Fehlerordnung. Warum ist der Mittelpunkt das \glqq beste Rechteck\grqq?
\tp{4}

\karo{6cm}

\clearpage
\textbf{(b)} Wir betrachten $I=\int_0^2 (x^2+1)\,dx$ mit St\"utzwerten $f(0)=1$, $f(1)=2$,
$f(2)=5$.
\tp{5}

\smallskip
Berechnen Sie den exakten Wert von $I$. Approximieren Sie $I$ anschlie\ss end mit der
\emph{Trapezregel} ($h=1$, zwei Teilintervalle) und der \emph{Simpson-Regel} ($h=1$). Vergleichen
Sie beide mit dem exakten Wert und erkl\"aren Sie, warum die Simpson-Regel hier exakt ist, die
Trapezregel aber nicht.

\karo{7.5cm}

\clearpage
\textbf{(c)} Bestimmen Sie das \emph{Lagrange}-Interpolationspolynom $p(x)$ durch die Punkte
$(0,1),(1,2),(2,5)$, vereinfachen Sie es und werten Sie $p(\tfrac12)$ aus.
\tp{3}

\smallskip
Was besagt das \emph{Runge-Ph\"anomen}, und welche Konsequenz hat es f\"ur die Wahl des
Polynomgrades bei der Quadratur (\glqq stop at Simpson\grqq)?

\karo{6.5cm}

\clearpage
\textbf{(d)} Erkl\"aren Sie die \emph{Monte-Carlo}-Integration mit dem Sch\"atzer
\[
   \hat I=\lvert\Omega\rvert\,\frac1N\sum_{i=1}^N f(X_i),\qquad
   \mathrm{SE}=\frac{\lvert\Omega\rvert\,\sigma_f}{\sqrt N}.
\]
Um welchen Faktor muss $N$ wachsen, um den Standardfehler zu \emph{halbieren}? Warum ist
Monte-Carlo im Hochdimensionalen vorteilhaft gegen\"uber gitterbasierter Quadratur
(Fluch der Dimensionen)? Inwiefern ist die Mittelpunktsregel ein Spezialfall von Monte-Carlo?
\tp{2}

\karo{6cm}

% =====================================================================
% AUFGABE 6 -- Kapitel 7
% =====================================================================
\clearpage
\aufg{6}{Modelltraining aus ersten Prinzipien}{10}

Wir trainieren ein lineares Regressionsmodell $\hat y(x)=w x+b$ mit dem mittleren quadratischen
Fehler $L(w,b)=\tfrac1n\sum_{i=1}^n (w x_i+b-y_i)^2$ per Gradientenabstieg. Der Datensatz der
Zielwerte sei $y=(2,\,4,\,6)$.

\smallskip
\textbf{(a)} Wir initialisieren mit $w=0$ und $b=\bar y$ (Mittelwert der Zielwerte). Zeigen Sie,
dass der erste Loss-Wert gerade die \emph{Varianz} der Zielwerte ist, und berechnen Sie ihn
konkret. Was w\"urde ein gedruckter erster Loss von $\approx 100$ statt des erwarteten Werts
bedeuten, was ein Wert von $\approx 0$?
\tp{4}

\karo{6cm}

\clearpage
\textbf{(b)} Eine Standardempfehlung ist, den Loss auf einer \emph{kleinen} Teilmenge mit einem
\emph{\"uberparametrisierten} Modell auf nahezu null zu treiben (overfit sanity check). Erkl\"aren
Sie, warum ein \emph{Scheitern} dieses Tests auf einen Fehler in Modell/Loss/Datenpipeline
hindeutet -- und \emph{nicht} im Optimierer. Welcher allgemeinen numerischen Praxis entspricht das?
\tp{3}

\karo{5.5cm}

\clearpage
\textbf{(c)} Ordnen Sie jede der folgenden Trainingspraktiken der zugrunde liegenden
\emph{numerischen Methode} aus der Vorlesung zu und begr\"unden Sie kurz:
\tp{3}
\begin{enumerate}\setlength{\itemsep}{1pt}
\item[(i)]   Mini-Batch-SGD (Gradient aus einer zuf\"alligen Teilmenge der Daten).
\item[(ii)]  Training mit mehreren zuf\"alligen Seeds, bester Lauf gewinnt.
\item[(iii)] Quantisierung des Modells f\"ur die Inferenz (niedrigere Bit-Pr\"azision).
\end{enumerate}
Geben Sie f\"ur (i) zus\"atzlich an, wie der Sch\"atzfehler des Gradienten mit der Batch-Gr\"o\ss e
$B$ skaliert.

\karo{6cm}

\end{document}
